% Generated by convert_myst_to_latex.py from chapter07/chapter07_sec03.md.
% Change the master and convert again instead of editing this file.
\RequirePackage[l2tabu, orthodox, abort]{nag}
\documentclass[paper=a4, fontsize=11pt, parskip=half]{scrartcl}

% exercises_preamble.tex
%
% Präambel für die Übungsblätter zum Selbststudium (Kapitel N.3 bzw. N.5) als
% PDF, übernommen aus der Präambel der CA-2-Unterlagen
% (internal/exams/exam_2026ws/ca02/preamble.tex). Die Studierenden üben damit
% das Format von CA 3: Aufgaben auf Papier, gelöst lokal am PC.
% Die .tex-Dateien erzeugt convert_myst_to_latex.py aus dem Master.

% ── Seitenlayout ─────────────────────────────────────────────────────────
\usepackage[a4paper, width=150mm, height=240mm]{geometry}

\usepackage{amsmath, amssymb}

\usepackage[no-math]{fontspec}
\usepackage{unicode-math}
\usepackage{libertinus}

\usepackage{setspace}
\usepackage{microtype}
\usepackage{textcomp}
\usepackage[main=english]{babel}

% ✩ für den Schwierigkeitsgrad (Libertinus hat das Zeichen nicht)
\usepackage{pifont}
\usepackage{newunicodechar}
\newunicodechar{✩}{\ding{73}}

% ── Kopf-/Fußzeilen ──────────────────────────────────────────────────────
\usepackage[automark]{scrlayer-scrpage}
\clearpairofpagestyles

\ifoot{S. Gramsch}
\cfoot{TH Mannheim -- Applied Numerical Mathematics -- WS 2026/27}
\ofoot{\pagemark}

% Keine Titelei: Die erzeugte .tex-Datei beginnt mit der Überschrift
% "N.M Exercises" als \section*, darunter "How to Practice" und "Exercises"
% als \subsection*.

% ── Listen ───────────────────────────────────────────────────────────────
\usepackage{tabularx}
\usepackage{booktabs}
\usepackage{enumitem}

% ── Farben ───────────────────────────────────────────────────────────────
\usepackage{xcolor}
\definecolor{itwm_blue_04}{HTML}{005A94}
\definecolor{frankfurtBlueLight}{HTML}{AECAE2}

\definecolor{my_darkgray}{HTML}{484949}
\definecolor{my_lightgray}{HTML}{F3F4F4}
\definecolor{my_orange}{HTML}{E87846}

% Farben der Admonition "tip" im Buch (MyST book-theme, heller Modus):
% Rand und Symbol --myst-color-success, Kopfzeile --myst-color-success-bg
\definecolor{tip_green}{HTML}{16A34A}
\definecolor{tip_bg}{HTML}{F0FDF4}

% ── Python-Code ──────────────────────────────────────────────────────────
\usepackage{listings}

% Zeilenumbrüche im Code mit Pfeil markieren, damit beim Abtippen klar ist,
% dass es eine einzige Python-Zeile ist
\lstdefinestyle{python}{
  language=Python,
  basicstyle=\ttfamily\small\setstretch{1},
  keywordstyle=\color{itwm_blue_04}\bfseries,
  stringstyle=\color{my_orange!80!black},
  commentstyle=\color{gray},
  showstringspaces=false,
  columns=fullflexible,
  keepspaces=true,
  upquote=true,
  breaklines=true,
  postbreak=\mbox{\textcolor{gray}{$\hookrightarrow$}\space},
  extendedchars=true,
}

\lstdefinestyle{ausgabe}{
  basicstyle=\ttfamily\footnotesize\setstretch{1},
  columns=fullflexible,
  keepspaces=true,
  upquote=true,
}

\lstdefinestyle{inline}{
  basicstyle=\ttfamily,
  showstringspaces=false,
  columns=fullflexible,
  keepspaces=true,
  upquote=true,
}

% Code im Fließtext: \py|derivative(values, dt)|
\newcommand{\py}{\lstinline[style=inline]}

% ── Kästen ───────────────────────────────────────────────────────────────
\usepackage[most]{tcolorbox}
\tcbuselibrary{listings}

%% Aufgabe ---------------------------------------------------------------
% \begin{aufgabe}{Exercise 7.1 (✩)}{approx.\ 7 min}: Bezeichnung links,
% Richtzeit rechts (in CA 2 und CA 3 stehen dort die Punkte).
% Aussehen wie die Admonition "tip" im Buch: grüner Rand links, hellgrüne
% Kopfzeile mit grünem Stift, weißer Inhalt, leichter Schatten.
% Eine Aufgabe wird nicht über zwei Seiten geteilt: Passt sie nicht mehr,
% beginnt sie auf der nächsten Seite. Nur für Aufgaben, die länger als eine
% Seite sind, setzt der Konverter die teilbare Variante aufgabelang ein.
\tcbset{
  aufgabenstil/.style 2 args={
    enhanced,
    boxrule=0pt,
    frame hidden,
    borderline west={1.5pt}{0pt}{tip_green},
    title={\textcolor{tip_green}{\ding{46}}\enspace #1 \hfill #2},
    coltitle=black,
    fonttitle=\bfseries,
    colbacktitle=tip_bg,
    colback=white,
    drop fuzzy shadow=black!25,
    arc=1mm,
    left=4mm, right=4mm,
    before skip=4mm, after skip=4mm,
  },
}
\newtcolorbox{aufgabe}[2]{aufgabenstil={#1}{#2}}
\newtcolorbox{aufgabelang}[2]{aufgabenstil={#1}{#2}, breakable}

%% Code und Ausgabe ------------------------------------------------------
% breakable: längerer vorgegebener Code darf über eine Seite hinausgehen,
% sonst bleibt vor ihm eine halbe Seite leer
\newtcblisting{code}{
    listing only,
    enhanced,
    breakable,
    boxrule=0pt,
    frame hidden,
    sharp corners,
    colback=my_lightgray,
    left=4pt, right=4pt, top=2pt, bottom=2pt,
    listing options={style=python},
}

\newtcblisting{ausgabe}{
    listing only,
    enhanced,
    boxrule=0pt,
    frame hidden,
    sharp corners,
    colback=white,
    borderline west={2pt}{0pt}{my_lightgray!70!black},
    left=6pt, right=4pt, top=0pt, bottom=0pt,
    listing options={style=ausgabe},
}

% ── Verweise ─────────────────────────────────────────────────────────────
\usepackage[hidelinks]{hyperref}


\begin{document}
\section*{7.3 Exercises}

\subsection*{How to Practice}

\begin{itemize}[noitemsep]
  \item Open a new Jupyter notebook on your own computer, switch off the
        internet connection and do not use AI tools.
  \item Type the given code yourself and pay attention to indentation,
        parentheses and quotation marks. A gray arrow
        \textcolor{gray}{$\hookrightarrow$} marks a line that is too long for
        the page. Type it as one line.
  \item Write your predictions and explanations in Markdown cells.
\end{itemize}

\subsection*{Exercises}

These exercises are meant for self-study at home and review the material of Chapters 7.1 and 7.2. Allow about 45 minutes for them. Before starting the exercises, run the following cell. It contains the function \py|derivative()| from Chapter 7.1 and the functions \py|motion_profile()| and \py|velocity()| of the smooth start from Chapter 7.2.

The difficulty level is given in the title of each exercise:

\begin{itemize}
  \item ✩ Comprehension: predict and explain code and outputs (approx. 7 min)
  \item ✩✩ Application: write your own code and interpret results (approx. 25 min)
\end{itemize}

\begin{code}
import numpy as np
import matplotlib.pyplot as plt

# function from Chapter 7.1
def derivative(values, dt):
    """Derivative of a measurement series with step size dt.

    Central difference quotient at the interior points, forward at the
    first point, backward at the last point.
    """
    n = len(values)
    result = np.zeros(n)
    result[0] = (values[1] - values[0]) / dt
    result[-1] = (values[-1] - values[-2]) / dt
    for k in range(1, n - 1):
        result[k] = (values[k + 1] - values[k - 1]) / (2 * dt)
    return result

# functions from Chapter 7.2
def motion_profile(t):
    """Distance in m for the smooth start, t in s from 0 to 2."""
    return 0.4 * (t - 2 / np.pi * np.sin(np.pi * t / 2))

def velocity(t):
    """Exact velocity in m/s, the derivative of motion_profile."""
    return 0.4 * (1 - np.cos(np.pi * t / 2))
\end{code}

\begin{aufgabe}{Exercise 7.1 (✩)}{approx.\ 7 min}
With \textbf{average speed enforcement} on highways, cameras photograph every vehicle at several checkpoints. From the distance between two checkpoints and the time between the two photos, the system computes the speed. The following code is given for a car that passes three checkpoints:

\begin{code}
import numpy as np

# location of the checkpoints in km and passing times in min
location = np.array([0.0, 6.0, 15.0])
time = np.array([0.0, 2.5, 7.0])

# speed on each section in km/h
v_section = np.zeros(len(location) - 1)
for k in range(len(location) - 1):
    v_section[k] = (location[k + 1] - location[k]) / (time[k + 1] - time[k]) * 60
print(v_section)

# speed over the entire route in km/h
v_total = (location[-1] - location[0]) / (time[-1] - time[0]) * 60
print(round(v_total, 1))
\end{code}

Write down your guess in a Markdown cell before you run the code.

\begin{enumerate}
  \item What do the two \py|print()| lines output? What is the factor \py|60| for?
  \item The speed limit is $130\,\text{km/h}$. Does the driver get a speeding ticket if the system evaluates both sections separately? And if it only evaluates the entire route?
  \item On the second section, the result is $120\,\text{km/h}$. Could the driver still have driven faster than $130\,\text{km/h}$ there? Explain your answer using the terms difference quotient and derivative.
  \item Run the code and check your predictions.
\end{enumerate}

\end{aufgabe}

\begin{aufgabe}{Exercise 7.2 (✩)}{approx.\ 7 min}
The AGV is recharging at its charging station. Every $10\,\text{min}$, the charging station logs the state of charge of the battery in percent. The following code is given:

\begin{code}
# state of charge in percent at 0, 10, 20, 30 and 40 min
dt = 10
state_of_charge = np.array([50.0, 65.0, 80.0, 90.0, 95.0])
print(derivative(state_of_charge, dt))
\end{code}

Write down your guess in a Markdown cell before you run the code.

\begin{enumerate}
  \item Which difference quotient does \py|derivative()| use for the first, the second and the last value? What does the \py|print()| line output, and in what unit?
  \item Which point in time does the last value fit best?
  \item What does the trend of the values tell us about the charging?
  \item Run the code and check your predictions.
\end{enumerate}

\end{aufgabe}

\begin{aufgabe}{Exercise 7.3 (✩)}{approx.\ 7 min}
The following code is given for the smooth start from Chapter 7.2:

\begin{code}
def backward(function, t, dt):
    """Backward difference quotient of function at the point t."""
    return (function(t) - function(t - dt)) / dt

t_0 = 1.0
exact = velocity(t_0)
for dt in [0.1, 0.05, 0.025]:
    deviation = backward(motion_profile, t_0, dt) - exact
    print(f'dt = {dt:5}: deviation {deviation:+.5f} m/s')
\end{code}

Write down your guess in a Markdown cell before you run the code.

\begin{enumerate}
  \item Expand $s(t - \Delta t)$ with the Taylor expansion from Chapter 7.2 up to the term with $\Delta t^2$. To do so, replace $\Delta t$ with $-\Delta t$ everywhere in the expansion of $s(t + \Delta t)$. Substitute the result into the backward difference quotient $(s(t) - s(t - \Delta t))/\Delta t$. Which error term is the largest one left?
  \item At $t = 1\,\text{s}$, the acceleration reaches its maximum of $s''(1) = 0.628\,\text{m/s}^2$. Which deviations, including their signs, do you predict for the three step sizes?
  \item By what factor does the deviation change from line to line?
  \item Run the code and check your predictions.
\end{enumerate}

\end{aufgabe}

\begin{aufgabe}{Exercise 7.4 (✩✩)}{approx.\ 25 min}
An AGV drives straight through the warehouse. At $t = 0.4\,\text{s}$, the safety scanner from Chapter 6.2 detects an obstacle in the protective field, and the AGV brakes as hard as it can. An unsecured steel wire mesh container stands on its load platform. With a coefficient of static friction of $\mu = 0.2$, the container slides as soon as the deceleration exceeds $\mu \cdot g = 0.2 \cdot 9.81\,\text{m/s}^2 \approx 1.96\,\text{m/s}^2$. During the drive, the AGV logged its position ten times per second:

\begin{code}
# step size in s and times 0.0, 0.1, ..., 1.5 s
dt = 0.1
t = np.linspace(0, 1.5, 16)
# distance traveled in m, rounded to millimeters
s = np.array([0.0, 0.1, 0.2, 0.3, 0.4, 0.488, 0.55, 0.588,
              0.6, 0.6, 0.6, 0.6, 0.6, 0.6, 0.6, 0.6])
\end{code}

\begin{enumerate}
  \item Use \py|derivative()| to compute the velocity. How fast was the AGV moving before it braked? How long is the braking distance from the start of braking at $t = 0.4\,\text{s}$ until the AGV stops?
  \item Compute the acceleration and plot velocity and acceleration in two subplots, one above the other. Draw the limit above which the container slides as a horizontal line, like the planned acceleration in Chapter 7.1. Determine the largest deceleration with \py|np.min()|. Does the container slide?
  \item Answer without code: the AGV brakes as hard as it can, abruptly and with constant deceleration. Why does the computed deceleration nevertheless rise gradually over several points in time and then drop again?
  \item With constant deceleration $a$, the braking distance is $v_0^2 / (2a)$. Check the formula with your results from tasks 1 and 2. The fleet manager is supposed to limit the deceleration to $1.96\,\text{m/s}^2$ so that the container does not slide. How long will the braking distance be then, and what does that mean for the protective field?
\end{enumerate}

Structure your code with IPO comments (input, processing, output).

\end{aufgabe}

\end{document}
